\documentclass[10pt,a4paper]{amsart}
\usepackage[margin=19mm]{geometry}
\usepackage{amsmath,amssymb,mathtools}
\usepackage[scr=rsfs,cal=euler]{mathalfa}
\usepackage{xfrac}
\usepackage[unicode,hidelinks,bookmarks=false]{hyperref}
\newcommand{\ie}{i.e.\ }
\newcommand{\cf}{cf.\ }
\newcommand{\resp}{respectively\ }
\newcommand{\Subsection}{\subsection}
\providecommand{\nobreakdash}{\nobreak\hbox{-}}
\setlength{\emergencystretch}{2em}
\multlinegap0pt
\usepackage{bm}
\usepackage{bbm}
\usepackage{dsfont}
\usepackage{xspace}
\usepackage{cancel}
\usepackage{mathtools}
\usepackage{amsthm}
\makeatletter
\@ifundefined{thm}{\newtheorem{thm}{Theorem}}{}
\@ifundefined{prop}{\newtheorem{prop}[thm]{Proposition}}{}
\@ifundefined{rmk}{\newtheorem{rmk}[thm]{Remark}}{}
\makeatother
\def\R{\mathbb R}
\def\Z{\mathbb Z}
\newcommand{\eps}{\varepsilon}
\newcommand{\ii}{{\mathbf i}}
\renewcommand\ge{\geq}
\renewcommand\le{\leq}
\title[Leaders in multi-type TASEP]{Leaders in multi-type TASEP}
\author{Alexei Borodin, Alexey Bufetov}
\date{Source: arXiv:2601.21412v1 (CC BY 4.0)}
\begin{document}
\maketitle

\noindent\emph{Source and licence.} Leaders in multi-type TASEP. Version arXiv:2601.21412v1 (CC BY 4.0), distributed under the Creative Commons Attribution 4.0 International licence, as recorded in the arXiv licence metadata for this exact version and shown on the abstract page https://arxiv.org/abs/2601.21412v1. Full rights notice: licenses/arxiv-2601.21412-CC-BY-4.0.txt.\par\smallskip

\noindent\textit{Editorial scope.} Counted unit: the Rmk whose source label is \texttt{rem:permut} in the article (source0.tex lines 1148--1156, source label \texttt{rem:permut}), delivered together with the article's own complete original proof of it (source0.tex lines 1158--1213, one \texttt{proof} environment of 56 lines). No mathematics is written, repaired or re-derived: statement and proof are the article's own text line for line. Required context retained uncounted: prop:2leaders-1 (lines 874--884); eq:prop2leaders1 (lines 877--882); eq:asymp2joint1 (lines 898--901); prop:2leaders-2 (lines 931--939); th:joint (lines 1118--1146); eq:density2a (lines 1133--1139). Every definition, display and label of those lines is retained unchanged. Omitted from this package and not redistributed: 1--873, 883--897, 902--930, 940--1117, 1140--1147, 1157--1157, 1214--1615. Nothing inside a retained excerpt is omitted or altered: every retained body is delivered exactly as the article's lines are written, so no omission receipt of any kind accompanies this package. Citations: the retained excerpts carry no citation command at all, so no bibliography entry of the article is transcribed and no reference is redistributed. Reference adaptation: none was needed and none was made; every reference inside the delivered unit resolves inside it, so no unresolved reference or citation is produced. Printed slips kept literal: no printed word, symbol, hypothesis or number of the counted statement or of its proof was altered. Layout and class: the article's own class is \texttt{amsart} with packages including amsfonts, amssymb, amsmath, epic, eepic, float, fullpage, longtable, color, siunitx, empheq, euscript, fontenc, mathtools; the wrapper is the lane's pinned \texttt{amsart} editorial wrapper at 10pt on A4, which restates the theorem-like environments the retained text needs and adds the article's own macro definitions that the retained text needs (\texttt{\textbackslash R}, \texttt{\textbackslash Z}, \texttt{\textbackslash eps}, \texttt{\textbackslash ge}, \texttt{\textbackslash ii}, \texttt{\textbackslash le}), with extra wrapper packages (bm, bbm, dsfont, xspace, cancel, mathtools). The delivered numbering is the unit's own. Assets: no retained span contains a figure environment, a graphics-inclusion command, a PDF-page inclusion, a table or a TikZ picture; no image file is copied into this package, the package declares no asset dependency and no image file is redistributed with it. Licence: version arXiv:2601.21412v1 of this article is distributed under the Creative Commons Attribution 4.0 International licence, as recorded in the arXiv licence metadata for this exact version and shown on the abstract page https://arxiv.org/abs/2601.21412v1. A full rights notice is delivered at licenses/arxiv-2601.21412-CC-BY-4.0.txt. Characters: every retained excerpt is pure ASCII, so the delivered engine, XeLaTeX with its default Latin Modern text font, reports no missing character.
	\begin{prop}
		\label{prop:2leaders-1}
		Let $k_3 >k_2 > k_1$ be arbitrary integers. We have the following equality:
		\begin{multline}
			\label{eq:prop2leaders1}
			\mathrm{Prob} \left( \mbox{the $k_1$-leader has type $\ge k_3$, } \mbox{the $(k_1, 2)$-leader has type $\ge k_2$} \right)
			\\ =
			\frac{1}{(2 \pi \ii)^2} \oint_{|v_2|=2} \oint_{|v_3|=2} \frac{(v_3-v_2)(v_2 v_3^2-1)(v_2^2 v_3-1) v_2^{k_1-k_2} v_3^{k_1-k_3}}{ v_2^2 v_3 (v_2 v_3-1) (v_2-1) (v_3-1)^3 } \\ \times \exp \left( t \left( \frac{1}{v_2 v_3} + v_2 + v_3 - 3 \right) \right) d v_2 d v_3.
		\end{multline}
		where the integration contours are counterclockwise.
	\end{prop}

		\begin{equation}
			\label{eq:asymp2joint1}
			\mathrm{Prob} \left( \mbox{the $k_1$-leader has type $\ge k_3$, } \mbox{the $(k_1, 2)$-leader has type $\ge k_2$} \right).
		\end{equation}

	\begin{prop}
		\label{prop:2leaders-2}
		Let $k_3 > k_2 >k_1$ be arbitrary integers. We have the following equality:
		\begin{multline*}
			\mathrm{Prob} \left( \mbox{the $k_1$-leader has type $\in [k_2;k_3)$,} \mbox{the $(k_1, 2)$-leader has type $\ge k_3$} \right) 
			\\ = \frac{1}{(2 \pi \ii)^2} \oint_{|v_2|=2} \oint_{|v_3|=2} \frac{(v_2 v_3^2-1)(v_2^2 v_3-1) \left( v_2^{k_1-k_2} v_3^{k_1-k_3} - v_2^{k_1-k_3} v_3^{k_1-k_2} \right) }{ v_2^2 v_3 (v_2 v_3-1) (v_2-1) (v_3-1)^2 } \\ \times \exp \left( t \left( \frac{1}{v_2 v_3} + v_2 + v_3 - 3 \right) \right) d v_2 d v_3.
		\end{multline*}
		
	\end{prop}

	\begin{thm}
		\label{th:joint}
		Let us fix $k_1 \in \Z$. Let $L_1 (t), L_2 (t) $ be the types of the $(k_1,1)$- and $(k_1,2)$-leaders, respectively, at time $t$. We have the following convergence in distribution:
		\begin{equation}
			\lim_{t \to \infty} \left( \frac{L_1(t) - k_1}{\sqrt{t}}, \frac{L_2(t) - k_1}{\sqrt{t}} \right) = G (a_2, a_3) d a_2 d a_3,
		\end{equation}
		where the limiting density $G (a_2, a_3)$ is given via
		$$
		G (a_2, a_3) :=
		\begin{cases}
			d_1 (a_2, a_3), \qquad \mbox{for $a_2 > a_3$,} \\
			d_2 (a_2,a_3), \qquad \mbox{for $a_2 \le a_3$,}
		\end{cases}
		$$
		and
		\begin{multline}
			\label{eq:density2a}
			d_1 (a_2,a_3) :=
			\frac{-1}{16 \pi} \left( \left( 4 \cdot \mathrm{erf} \left( \frac{\sqrt{3}}{6} (a_2 +a_3) \right) -4 \right) \exp \left( \frac{(a_2-a_3)^2}{4} \right) \sqrt{\pi} (a_2 - a_3)
			\right. \\ \left. + 2 \sqrt{3} \left( a_2^2 -6 \right) \exp \left( \frac{-a_2^2+a_2 a_3 - a_3^2}{3} \right)
			\right. \\ \left. + \sqrt{\pi} \exp \left( \frac{-a_2^2}{4} \right) \left(  \mathrm{erf} \left( \frac{\sqrt{3}}{6} (a_2 - 2 a_3) \right) +1 \right) \left( - 4 a_2 +a_2^3 - 2 a_2^2 a_3 + 4 a_3 \right) \right. ,
		\end{multline}
		\begin{multline*}
			d_2 (a_2,a_3) := \frac{1}{4 \sqrt{\pi}} \left( - \exp \left( -\frac14 (a_2-a_3)^2 \right) (a_2 - a_3) \left( \mathrm{erf} \left( \frac{\sqrt{3}}{6} (a_2+a_3) \right) -1 \right)
			\right. \\ \left. + a_3 \exp \left( \frac{-a_3^2}{4} \right) \mathrm{erf} \left( \frac{\sqrt{3}}{6} (2 a_2-a_3) \right) +a_2 \exp \left( \frac{-a_2^2}{4} \right) \mathrm{erf} \left( \frac{\sqrt{3}}{6} ( a_2- 2 a_3) \right)
			\right. \\ \left. - a_3 \exp \left( \frac{-a_3^2}{4} \right) + a_2 \exp \left( \frac{-a_2^2}{4} \right) \right).
		\end{multline*}
		
	\end{thm}

	\begin{rmk}
		\label{rem:permut}
		By a direct integration, it follows from Theorem \ref{th:joint} that
		$$
		\lim_{t \to \infty} \mathrm{Prob} \left( L_1(t) > L_2 (t) \right) = 
		\lim_{t \to \infty} \mathrm{Prob} \left( L_1(t) < L_2 (t) \right) = \frac12. 
		$$
		We are unaware of a conceptual explanation or a shorter proof of this equality; it would be interesting to understand the origin of this fact better. A naive generalization for the ordering of types of more than two leaders --- that they always form a uniformly random permutation --- is false. 
	\end{rmk}

	\begin{proof}
		
		\textbf{First step}. 
		
		\smallskip
		
		Choose the sequences $k_2 = k_2(t)$, $k_3 = k_3 (t)$ so that $k_1 - k_2 = - \lfloor a_2 \sqrt{t} \rfloor$, $k_1 - k_3 =- \lfloor a_3 \sqrt{t} \rfloor$ (and $k_1$ is fixed), and
		let us perform the asymptotic analysis of the expression given by Proposition \ref{prop:2leaders-1} via the steepest descent method, as $t \to \infty$. 
		
		The integrand of \eqref{eq:prop2leaders1} depends on $t$ only through the function $\exp ( t \Phi (v_2,v_3)) v_2^{-a_2 \sqrt{t}} v_3^{-a_3 \sqrt{t}}$, where
		$$
		\Phi (v_2,v_3) := \frac{1}{v_2 v_3} + v_2 + v_3 - 3.
		$$
		A direct computation shows that for $|v_2|=1$ and $|v_3|=1$ we have $\Re \left( \Phi (v_2,v_3) \right) \le 0$, and the equality is reached for $v_2=v_3=1$ only. Also note that $(v_2,v_3)=(1,1)$ is a critical point of $\Phi (v_2,v_3)$. 
		
		Now, let us deform the contours in the expression from \eqref{eq:prop2leaders1} to circles $|v_2|=|v_3|=1 +1/ t^{10}$ (we do not choose exactly $1$ in order to avoid the pole at $v_2 v_3 =1$). Standard for the steepest descent method estimates show that for these contours the leading contribution to the integral comes from an order $t^{-1/2}$ neighborhood of the critical point $(v_2,v_3)=(1,1)$. 
		We make a change of variables $v_3 = 1+ \frac{u_3}{\sqrt{t}}$, $v_2 = 1+ \frac{u_2}{\sqrt{t}}$ in order to compute the integral in this neighborhood. We have
		$$
		\Phi (v_2,v_3) = 2 + \frac{u_3+u_2}{\sqrt{t}} + \frac{1}{1+\frac{u_3+u_2}{\sqrt{t}} + \frac{u_3 u_2}{t} + o(t^{-1}) } -3 = \frac{u_3^2+u_3 u_2 +u_2^2}{t} + o( t^{-1}).
		$$
		Plugging these expressions into the integral, we obtain that in the limit $t \to \infty$ the probability from \eqref{eq:asymp2joint1} is equal to
		\begin{multline*}
			\frac{1}{(2 \pi \ii)^2} \iint \frac{(u_3-u_2)(u_2 + 2 u_3)(2 u_2 + u_3)}{ (u_2 + u_3) u_2 u_3^3 } \\ \times \exp \left( u_2^2 + u_3 u_2 +u_3^2 - a_2 u_2 - a_3 u_3) \right) d u_2 d u_3 + O \left( \frac{1}{\sqrt{t}} \right),
		\end{multline*}
		where the integration contours are $\eps+ \ii \R$ lines, for $0 < \eps \ll 1$. Note that $t^{-1/2}$ appeared thrice in the numerator, five times in the denominator, and two times from Jacobian of the change of variables, which led to the cancellation of this factor in the leading term of the asymptotics.
		After the change of variables $u_2 = \ii x_2$, $u_3 = \ii x_3$, we obtain the following expression:
		\begin{multline*}
			Y(a_2,a_3) := \frac{1}{(2 \pi \ii)^2} \iint \frac{(x_3-x_2)(x_2 + 2 x_3)(2 x_2 + x_3)}{ (x_2 + x_3) x_2 x_3^3 } \\ \times \exp \left( - x_2^2 - x_3 x_2 - x_3^2 - \ii a_2 x_2 - \ii a_3 x_3 \right) d x_2 d x_3,
		\end{multline*}
		where the integration contours are real axes slightly shifted upwards.
		
		The limiting density for the vector 
		$$
		\left( \frac{L_1(t) - k_1}{\sqrt{t}}, \frac{L_2(t) - k_1}{\sqrt{t}} \right)
		$$
		in the $L_2 (t) < L_1(t)$ region is given by
		$$
		\partial_{a_2} \partial_{a_3} Y(a_2,a_3).
		$$
		We claim that this expression is equal to $d_1(a_2,a_3)$ given by \eqref{eq:density2a}. One way to compute the integral is the following: by further differentiation of the integrand with respect to parameters $a_2$ and $a_3$, one can get rid of the denominator and compute the integral in $x_2, x_3$ explicitly. After this, one can take necessary antiderivatives in $a_2, a_3$. We leave details of this computation to math software.
		
		\smallskip
		
		\textbf{Second Step.}
		
		\smallskip
		
		In order to address the $L_2 (t) > L_1(t)$ region, we need to perform asymptotic analysis of the observable from Proposition \ref{prop:2leaders-2}. Performing exactly the same steps as in the first case, we obtain that the leading term of the asymptotics is given by the integral
		\begin{multline*}
			\hat Y (a_2,a_3) := \frac{1}{(2 \pi \ii)^2} \oint \oint \frac{(x_3-x_2)^2 (x_2 + 2 x_3)(2 x_2 + x_3)}{ (x_2 + x_3) x_3^3 } \\ \times \exp \left( - x_2^2 - x_3 x_2 - x_3^2 - \ii a_2 x_2 - \ii a_3 x_3 \right) d x_2 d x_3,
		\end{multline*}
		where again the integration contours are real axes slightly shifted upwards. 
		
		Evaluating this integral in a similar fashion, we arrive at the statement of the proposition. 
		
	\end{proof}

\end{document}
